\section{Constructions used in the group-theoretic arguments}\label{sec:constructions}

\subsection{The coefficient field and the base coordinate pair}
Set
\[
 K=\Q[a]/(a^2-a+1),\quad b=1-a,\quad ab=1,\quad b^2=-a.
\]
The integrally normalized Chebyshev polynomials are defined recursively by
$D_0=2$, $D_1=t$, $D_{n+1}=tD_n-D_{n-1}$. In characteristic zero,
$D_n(t)=2T_n(t/2)$. In particular,
\[
 D_3=t^3-3t,\quad D_4=t^4-4t^2+2,\quad D_5=t^5-5t^3+5t,
\quad D_7=t^7-7t^5+14t^3-7t.
\]
We consider the specific parametrizations
\[
 \gamma_T=(D_3,D_4-2,t^5-10t),\qquad
 \gamma_H=(D_3,D_5,D_7).
\]
The coefficient of $t$ in the second component of $\gamma_H$ is $5$.
The parametrization with coefficient $4$ is not covered by this result.

The affine map $M=(x,y+2,z-5x)$ sends $\gamma_T$ to
$\gamma_0=(D_3,D_4,D_5)$. In the normalized coordinates, introduce
\begin{align*}
 u&=axy+z,&v&=-y,&m&=bx^5,&s&=bx,\\
 \ell&=-1+ax^2-ax^4,&
 g&=-2x^5+(6-2a)x^3-2bx,\\
 P_x(U)&=sU^2+\ell U+g,&H&=P_x(u)+mv.
\end{align*}
The polynomial $P_x$ induces an automorphism of $K[x,T]/(x^5)$.
We use the following representative of its inverse:
\begin{align}
 F_x(T)={}&-T+x(bT^2+2a-2)+x^2(2aT^3-5aT)\notag\\
 &+x^3(-5T^4+15T^2-2a)\notag\\
 &+x^4\bigl(14bT^5+50(a-1)T^3+(29-24a)T\bigr).
 \label{eq:inverse-mod-five}
\end{align}
Direct computation gives both congruences
\[
 P_x(F_x(T))\equiv T\pmod{x^5},\qquad
 F_x(P_x(U))\equiv U\pmod{x^5}.
\]
Consequently,
\[
 W=\frac{u-F_x(H)}m,\qquad
 \chi_x(T)=\frac{T-P_x(F_x(T))}m
\]
are polynomials. For $A_0=(x,H,W)$, the inverse map is given by
\begin{align}
 U&=F_x(T)+mV,\notag\\
 V_0&=\chi_x(T)-(\ell+2sF_x(T))V-smV^2,\notag\\
 A_0^{-1}(x,T,V)&=(x,-V_0,U+axV_0).\label{eq:base-inverse}
\end{align}
Indeed, expanding the square gives $P_x(U)+mV_0=T$ and
$(U-F_x(T))/m=V$. For the composition in the other direction, one first
recovers $u,v$ and then $y,z$. This is a polynomial proof of invertibility;
constancy of the Jacobian alone is not used as a criterion.

Substitution into the curve gives
\[
 H(D_3,D_4,D_5)=t,\qquad W(D_3,D_4,D_5)=8-14t^2.
\]
Thus, with $h_T=H\circ M$ and $w_T=W\circ M$, we obtain
\begin{equation}\label{eq:trefoil-rectifier}
 \Phi_T=(h_T,\ x-h_T^3+3h_T,\ w_T-8+14h_T^2),
 \qquad \Phi_T\gamma_T=(t,0,0).
\end{equation}
Its inverse is constructed from \eqref{eq:base-inverse}, $M^{-1}$ and the
inverses of the outer shears. The component degrees of this map are
$(5,15,24)$, and its Jacobian is $1$. No optimality of these degrees is claimed.
The original implementation is in the repository's \texttt{verify.py};
the new verification code separately reproduces the compact certificates.

\subsection{Transfer with control of the forward and inverse maps}
The following mechanism is related to residual coordinates
\cite{BD1993,EV2001,Lewis2024}; we record the formulas needed for
reproducibility.

\begin{proposition}[A sufficient transfer criterion]\label{prop:transfer-app2}
Let $A=(x,H,W)$ and $A^{-1}=(x,F,G)$ be mutually inverse polynomial
maps over $k$. Let $0\ne m\in k[x]$, $p\in k[x,y]$, and
\[
 \beta=(x,p(x,y)+m(x)z,y).
\]
If $c(y)=\overline{H(x,p(x,y),y)}$ defines an automorphism of
$(k[x]/(m))[y]$, then $H\circ\beta$ extends to a relative
polynomial automorphism.
\end{proposition}
\begin{proof}
Let $\psi$ be the inverse of $c$ in $B[T]$, where $B=k[x]/(m)$.
Choose a lift of the polynomial
\[
 \varphi(x,T)\equiv W(x,p(x,\psi(T)),\psi(T))\pmod m.
\]
Then set
\begin{equation}\label{eq:transfer-forward}
 A_{\rm new}=
 \left(x,H\circ\beta,
 \frac{W\circ\beta-\varphi(x,H\circ\beta)}m\right).
\end{equation}
The numerator of the last component vanishes modulo $m$.
For input $(x,T,V)$, put
$S=mV+\varphi(x,T)$, $Y=F(x,T,S)$, $Z=G(x,T,S)$. The inverse map is
\begin{equation}\label{eq:transfer-inverse}
 A_{\rm new}^{-1}(x,T,V)=\left(x,Z,\frac{Y-p(x,Z)}m\right).
\end{equation}
In $B[T]$, invertibility of $A$ yields
$Y=p(x,\psi(T))$, $Z=\psi(T)$, so the division is again exact.
After inverting $m$, these formulas are compositions of explicitly
invertible maps. Injectivity of localization carries both identity
compositions back to $k[x,y,z]$, including all fibres with $m=0$.
\end{proof}

The condition on $c=c_0+c_1y+\sum_{j\ge2}c_jy^j$ is equivalent to
$c_1$ being a unit and all $c_j$, $j\ge2$, being nilpotent in $B$.
If $J=(c_j:j\ge2)$ and $J^e=0$, the inverse is obtained by the finite iteration
\[
 \psi_0=c_1^{-1}(T-c_0),\qquad
 \psi_{i+1}=c_1^{-1}\left(T-c_0-\sum_{j\ge2}c_j\psi_i^j\right).
\]
Differences between successive iterates lie in successively higher powers
of $J$; after $e$ steps a fixed point is reached. Thus the computations
must be performed modulo the full $m$, rather than only its radical.

\subsection{The harmonic model and the family $C_N$}\label{subsec:family-app2}
The identity
\[
 D_3D_7-D_5^2+2=D_4
\]
shows that $\beta_1=(x,xz-y^2+2,y)$ sends $\gamma_H$ to
$\gamma_0$. Since $H(0,Y,Z)=-Z$, the transfer centre is $-y$
modulo $x$, and Proposition~\ref{prop:transfer-app2} applies.
One may take the correction
\[
 \varphi_1(T)=42(a-1)T^6+175(1-a)T^4+(147a-161)T^2+8.
\]
Putting $H_1=H\circ\beta_1$ and
$W_1=(W\circ\beta_1-\varphi_1(H_1))/x$, we have
\[
 W_1\gamma_H=b(42t^3-49t),\quad
 \Phi_H=(H_1,x-H_1^3+3H_1,W_1-b(42H_1^3-49H_1)).
\]
The component degrees are $(7,21,41)$. The inverse is given by
\eqref{eq:transfer-inverse} and the inverse outer shears.

Now let $N\ge1$. In $\Z[x,y]/(x^N)$, define
\[
 \nu_0=-y,\qquad \nu_{i+1}=-y+x(\nu_i^2-2),\quad 0\le i<N,
 \qquad p_N\equiv2+x\nu_N-\nu_N^2\pmod{x^N},
\]
choosing the remainder $p_N$ of degree less than $N$ in $x$.
We have $p_N\equiv2-y^2\pmod x$. Set
\begin{equation}\label{eq:family-CN}
 C_N(t)=\left(D_3,D_5,
 \frac{D_4-p_N(D_3,D_5)}{D_3^N}\right),\qquad
 \beta_N=(x,p_N+x^Nz,y).
\end{equation}
The division in this parametrization is exact. Indeed,
$D_5=D_3(t^2-2)-t$ and $D_4=D_3t-t^2+2$.
At each step, the iteration error $\nu_i(D_3,D_5)-t$ acquires an
additional factor $D_3$. Hence
$\nu_N(D_3,D_5)\equiv t\pmod{D_3^N}$, which gives the required divisibility.

The centre $H(x,p_N,y)$ equals $-y$ modulo $x$. Therefore, modulo $x^N$,
its linear coefficient is a unit and all higher coefficients are nilpotent.
Proposition~\ref{prop:transfer-app2} constructs $A_N=(x,H_N,W_N)$ with
$H_N=H\circ\beta_N$ and $H_NC_N=t$. With $w_N=W_NC_N$, we obtain
\begin{equation}\label{eq:family-rectifier}
 \Phi_N=(H_N,x-D_3(H_N),W_N-w_N(H_N)),\qquad
 \Phi_NC_N=(t,0,0).
\end{equation}
Thus $C_N$ does indeed define a closed embedding.
For $N=1$, one may choose the transfer above and obtain
$C_1=\gamma_H$, $\Phi_1=\Phi_H$. Different representatives of the
correction may change the companion by a tame shear; all subsequent
statements concern maps constructed in this way.

\subsection{Integrality of the formulas}
All base coefficients lie in
$\mathcal O=\Z[a]/(a^2-a+1)$. The congruences for $P_x,F_x$ already hold
over $\mathcal O$, and $b^{-1}=a$. Division by $bx^5$ consists of
multiplication by $a$ and exact division by the monomial $x^5$.
The recursions for $p_N$ and for the inverse of the centre introduce no
numerical denominators: the linear coefficient of the centre is $-1$
modulo $x$ and is inverted by a finite geometric series modulo $x^N$.
Consequently, the polynomial maps and their inverses specialize to any field
containing a root of $a^2-a+1$, including characteristics $2$ and $3$.
This is a statement about the formulas; the wildness proofs below apply in
characteristic zero. Degrees may decrease after specialization.
For $\mathbb F_q$, a root exists if and only if
$\operatorname{char}\mathbb F_q=3$ or $q\equiv1\pmod3$:
in characteristic different from $3$, the element $-a$ is a nontrivial
cube root of unity. The absence of a root does not rule out other rectifiers.
